\section{Misc}

\subsection{Big Integer}

\subsubsection*{Utilitários}

\begin{lstlisting}
string stripLeadingZeros (const string &number) {
    int firstNonZero = number.find_first_not_of('0');
    if (firstNonZero == string::npos) return "0";
    return number.substr(firstNonZero);
}
\end{lstlisting}

\subsubsection*{Soma}

Assume que ambos os números são positivos.

\textbf{Complexidade:} $O(n)$.

\begin{lstlisting}
string addBig (const string &first, const string &second) {
    string result;
    int carry = 0;
    int i = first.size() - 1;
    int j = second.size() - 1;
    while (i >= 0 || j >= 0 || carry) {
        int digitSum = carry;
        if (i >= 0) digitSum += first[i--] - '0';
        if (j >= 0) digitSum += second[j--] - '0';
        result.push_back('0' + digitSum % 10);
        carry = digitSum / 10;
    }
    reverse(all(result));
    return stripLeadingZeros(result);
}
\end{lstlisting}

\subsubsection*{Subtração}

Assume que $first \geq second$.

\textbf{Complexidade:} $O(n)$.

\begin{lstlisting}
string subtractBig (const string &first, const string &second) {
    string result;
    int borrow = 0;
    int i = first.size() - 1;
    int j = second.size() - 1;
    while (i >= 0) {
        int digit = (first[i] - '0') - borrow;
        if (j >= 0)
            digit -= second[j--] - '0';
        if (digit < 0) {
            digit += 10;
            borrow = 1;
        } else {
            borrow = 0;
        }
        result.push_back('0' + digit);
        i--;
    }
    reverse(all(result));
    return stripLeadingZeros(result);
}
\end{lstlisting}

\subsubsection*{Multiplicação por Inteiro}

\textbf{Complexidade:} $O(n)$.

\begin{lstlisting}
string multiplyByInt (const string &number, int multiplier) {
    string result;
    ll carry = 0;
    for (int i = number.size() - 1; i >= 0; i--) {
        ll product = (number[i] - '0') * 1LL * multiplier;
        product += carry;
        result.push_back('0' + product % 10);
        carry = product / 10;
    }
    while (carry > 0) {
        result.push_back('0' + carry % 10);
        carry /= 10;
    }
    reverse(all(result));
    return stripLeadingZeros(result);
}
\end{lstlisting}

\subsubsection*{Divisão por Inteiro}

\textbf{Complexidade:} $O(n)$.

\begin{lstlisting}
pair<string, int> divideByInt (const string &number, int divisor) {
    string quotient;
    ll remainder = 0;
    for (char digit : number) {
        remainder = remainder * 10 + (digit - '0');
        quotient.push_back('0' + remainder / divisor);
        remainder %= divisor;
    }
    return {stripLeadingZeros(quotient), remainder};
}
\end{lstlisting}

\subsubsection*{Módulo por Inteiro}

\textbf{Complexidade:} $O(n)$.

\begin{lstlisting}
int modString (const string &number, int mod) {
    ll remainder = 0;
    for (char digit : number) 
        remainder = (remainder * 10 + (digit - '0')) % mod;
    return remainder;
}
\end{lstlisting}


\subsection{Contagem de Inversões}

\textbf{Complexidade:} $O(n\log n)$.

\begin{lstlisting}
ll count_inversions(vector<int>& a) {
    // Compressao de Coordenadas
    vector<int> vals = a;
    sort(all(vals));
    vals.erase(unique(all(vals)), vals.end());
    SegTree st(sz(vals)); // Segtree para soma e update = 1
    ll inv = 0;

    for (int x : a) {
        int rank = lower_bound(all(vals), x) - vals.begin();
        inv += st.query(rank + 1, vals.size() - 1); 
        st.update(rank, 1);
    }
    return inv;
}
\end{lstlisting}


\subsection{Compressão de Coordenadas}

Mapeia valores grandes para um intervalo compacto $[1, k]$ preservando a ordem relativa em $O(n \log n)$.

\textbf{Complexidade:} $O(n + m + \text{matches})$.

\begin{lstlisting}
vector<int> vals = a; // vetor original
sort(all(vals));
vals.erase(unique(all(vals)), vals.end());

auto get_id = [&](int x) {
    return lower_bound(all(vals), x) - vals.begin() + 1; // 1-indexed
};
// Para cada elemento a[i], use get_id(a[i])
\end{lstlisting}


\subsection{Algoritmo de Mo (Range Queries Offline)}

\textbf{Complexidade:} $O((n + q) \sqrt{n})$.

\begin{lstlisting}
const int BLOCK = 450; // ~ sqrt(N)
struct Query {
    int l, r, id;
    bool operator<(const Query& other) const {
        int b1 = l / BLOCK, b2 = other.l / BLOCK;
        if (b1 != b2) return b1 < b2;
        return (b1 & 1) ? r < other.r : r > other.r; // Otimizacao de zigue-zague
    }
};
void add(int idx) { /* Adiciona o elemento na estrutura de dados */ }
void remove(int idx) { /* Remove o elemento da estrutura de dados */ }
long long get_answer() { /* Retorna a resposta atual */ return 0; }
vector<ll solve_mo(vector<Query>& queries, const vector<int>& a) {
    int q = queries.size();
    vector<ll ans(q);
    sort(queries.begin(), queries.end());
    int cur_l = 0, cur_r = -1;
    for (auto& qry : queries) {
        while (cur_l > qry.l) add(--cur_l);
        while (cur_r < qry.r) add(++cur_r);
        while (cur_l < qry.l) remove(cur_l++);
        while (cur_r > qry.r) remove(cur_r--);
        ans[qry.id] = get_answer();
    }
    return ans;
}
\end{lstlisting}


\subsection{Meet-in-the-Middle}

\textbf{Complexidade:} $O(2^{n/2} \log(2^{n/2}))$.

\begin{lstlisting}
// Exemplo generico para subconjuntos (Subset Sum / Knapsack)
vector<ll generate_sums(const vector<ll& v) {
    int n = v.size();
    vector<ll sums;
    for (int mask = 0; mask < (1 << n); ++mask) {
        long long s = 0;
        for (int i = 0; i < n; ++i) {
            if (mask & (1 << i)) s += v[i];
        }
        sums.push_back(s);
    }
    sort(all(sums));
    return sums;
}
\end{lstlisting}